\chapter{Replication Exercises}
\label{app:replication}

These exercises guide the reader through replicating key doxonomic analyses using real data.

\section*{Exercise 1: Epistemic Herfindahl Index for Cable / Broadcast News}

\textbf{Goal}: Compute $\ehhi$ for a national television-news market. The US cable-news worked example is given below; the same methodology applies to UK broadcast news (BBC, Sky, ITV, GB News), Brazilian open TV (Globo, Record, SBT, Band), Indian news TV (Republic, Times Now, NDTV, Aaj Tak, India Today), or any other partitioned national market.

\textbf{Data}: Nielsen audience-share data (available from Pew Research Center's annual ``State of the News Media'' reports for the US; BARB for the UK; Kantar IBOPE Media for Brazil; BARC for India; ATM-style audience services in EU member states).

\textbf{Steps}:
\begin{enumerate}[nosep]
  \item Obtain average primetime viewership for the major channels in the chosen market (e.g., in the US: CNN, Fox News, MSNBC, and other cable news channels; comparable channel lists for the UK, Brazil, India, etc.).
  \item Compute market shares $s_k = \text{viewership}_k / \text{total viewership}$.
  \item Compute $\ehhi = \sum_k s_k^2$.
  \item Interpret: is the cable / broadcast news market concentrated by doxonomic standards ($\ehhi > 0.25$)?
  \item Repeat for 2005, 2010, 2015, 2020. Plot the trend.
\end{enumerate}

\section*{Exercise 2: Donor Network of Free-Market Think Tanks}

\textbf{Goal}: Construct a donor-institution network for a national free-market think-tank ecosystem. The US case is given below; the same methodology applies to the UK (IEA, Adam Smith Institute, Centre for Policy Studies, Institute for Free Trade, TaxPayers' Alliance), Brazil (Instituto Millenium, Instituto Mises Brasil, Instituto Liberal), India (CCS, Liberty Institute, Takshashila), and other Atlas-Network national ecosystems.

\textbf{Data}: IRS 990 filings via ProPublica Nonprofit Explorer; InfluenceWatch profiles. (UK Charity Commission filings; Brazilian and Indian equivalent registries where available.)

\textbf{Steps}:
\begin{enumerate}[nosep]
  \item Select 10 major think tanks for the chosen market (in the US: Heritage Foundation, Cato Institute, AEI, Hoover, Manhattan Institute, etc.; substitute the national-equivalent list for non-US replications).
  \item For each, identify the top 5 foundation donors from 990 filings.
  \item Construct a bipartite network (donors $\times$ institutions) with edge weights = grant amounts.
  \item Compute eigenvector centrality for all nodes.
  \item Identify narrative clusters using modularity maximization.
\end{enumerate}

\section*{Exercise 3: Semantic Drift of ``Freedom''}

\textbf{Goal}: Measure how the word ``freedom'' has changed meaning in American political discourse.

\textbf{Data}: Google Ngrams or a corpus of presidential speeches.

\textbf{Steps}:
\begin{enumerate}[nosep]
  \item Train word embeddings (Word2Vec or GloVe) on decade-specific subcorpora.
  \item Extract the embedding vector for ``freedom'' in each decade.
  \item Compute cosine distance between consecutive decades.
  \item Identify the decade with maximum drift.
  \item Examine nearest neighbors of ``freedom'' in each period to interpret the drift.
\end{enumerate}

\section*{Exercise 4: Meritocratic Belief Across Countries}

\textbf{Goal}: Compare meritocratic belief across 20 countries.

\textbf{Data}: World Values Survey (items on effort, success, fairness).

\textbf{Steps}:
\begin{enumerate}[nosep]
  \item Select relevant items from WVS Wave 7.
  \item Construct a meritocratic belief index (factor analysis or simple sum).
  \item Plot country-level means against Gini coefficients.
  \item Test the hypothesis: does higher inequality correlate with stronger meritocratic belief (legitimation function)?
\end{enumerate}

\section*{Exercise 5: Narrative Entropy of Climate Discourse}

\textbf{Goal}: Measure narrative diversity in media coverage of climate change over time.

\textbf{Data}: Media Cloud or GDELT.

\textbf{Steps}:
\begin{enumerate}[nosep]
  \item Collect climate-related articles from 2000--2024.
  \item Apply LDA topic modeling with $K = 10$ topics.
  \item Compute narrative entropy $H(t) = -\sum_k p_k(t) \log p_k(t)$ for each year.
  \item Plot $H(t)$ over time. Identify periods of low entropy (narrative capture).
  \item Correlate entropy changes with known doxonomic events (e.g., fossil fuel campaign spending).
\end{enumerate}
